\begin{helper}
For every \(\eta>0\) and every \(\varepsilon_1>0\), there is \(n_0\)
such that for every \(n>n_0\), if \(L=\log n\), then
\[
\frac{2}{L}\le\frac{\varepsilon_1}{2}
\]
and
\[
\frac{1}{(1+\eta)n^{1/4}\sqrt L}\le\frac{\varepsilon_1}{2}.
\]
\end{helper}
\begin{proof}
Since \(\log n\to\infty\), enlarge \(n_0\) so that
\[
L=\log n\ge \frac{4}{\varepsilon_1}
\]
for all \(n>n_0\).  Multiplying by the positive number
\(\varepsilon_1/2\) gives
\[
\frac{\varepsilon_1}{2}L\ge 2,
\]
and hence \(2/L\le \varepsilon_1/2\).

Also enlarge \(n_0\) so that \(L\ge 1\), \(n\ge 1\), and
\[
n^{1/4}\ge \frac{2}{\varepsilon_1}.
\]
The last condition is possible because \(n^{1/4}\to\infty\).  Since
\(\eta>0\), we have \(1+\eta\ge1\), and \(L\ge1\) gives
\(\sqrt L\ge1\).  Therefore
\[
(1+\eta)n^{1/4}\sqrt L\ge n^{1/4}\ge \frac{2}{\varepsilon_1}.
\]
Taking reciprocals of positive quantities yields
\[
\frac{1}{(1+\eta)n^{1/4}\sqrt L}
  \le \left(\frac{2}{\varepsilon_1}\right)^{-1}
  =\frac{\varepsilon_1}{2}.
\]
\end{proof}
